\documentclass[12pt]{article}
%Gummi|065|=)
\usepackage{amsmath, amsfonts, amssymb}
\usepackage[margin=0.5in]{geometry}
\usepackage{xcolor}
\usepackage{graphicx}
%\usepackage{graphicx}
\newcommand{\off}[1]{}
\DeclareMathSizes{20}{30}{20}{18}
\newcommand{\myhrule}{}

\newcommand{\two }{\sqrt[3]{2}}
\newcommand{\four}{\sqrt[3]{4}}

\newcommand{\dash}{
\begin{tikzpicture}[scale=1]
\draw (0,0)--(19,0);
\end{tikzpicture}
}

\newcommand{\sq}[3]{
\node at (#1+0.5,#2+0.5) {#3};
\draw (#1+0,#2+0)--(#1+1,#2+0)--(#1+1,#2+1)--(#1+0,#2+1)--cycle;
}

\usepackage{tikz}

\title{\textbf{Proposal: Factorial}}
\author{John D Mangual}
\date{}
\begin{document}

\fontfamily{qag}\selectfont \fontsize{20}{25}\selectfont

\maketitle

\noindent Some clever person turned the theory domino tilings into a fundamental object of mathematics and of nature.  For a long time there were really two 





\newpage

\fontfamily{qag}\selectfont \fontsize{12}{10}\selectfont

\begin{thebibliography}{}

\item ...


\end{thebibliography} 

\noindent The most amazing typo ever: \\

\begin{tikzpicture} [scale=0.5]

\foreach \a in {0,...,5}{
    \draw (\a,5-\a)--(1, \a + 1);
}

\end{tikzpicture}

\begin{verbatim}
\begin{tikzpicture} [scale=0.5]

\foreach \a in {0,...,5}{

    \draw (\a, 5-\a)--(1, \a + 1);

}

\end{tikzpicture}
\end{verbatim}

\end{document}